17 questions in 45 minutes. On the projection screen in front of you, a
series of images will be projected. Look at each image and answer the
following questions.

\begin{description}
\item[(1.1)] (Image 1, 2 minutes) Identify the drawn lines on the image.
  [Answer table: for line numbers 1, 2, 3, tick one of --- Celestial
  Equator, Ecliptic, Local Meridian, Local Vertical, Zero RA, Galactic
  Equator.] \hfill[6.0 pt]
\item[(1.2)] (Image 2, 2 minutes) This image was taken from the town of Fada
  in Chad. What is the latitude of this place? Select only one choice:
  \begin{description}
  \item[A.] $10^\circ$
  \item[B.] $17^\circ$
  \item[C.] $22^\circ$
  \item[D.] $33^\circ$
  \end{description}
  \hfill[3.0 pt]
\item[(1.3)] (Image 3, 3 minutes) Two of the following stars are transiting
  the meridian in the image. Tick the boxes next to those stars:
  Arcturus, Mizar, Spica, Regulus. \hfill[4.0 pt]
\item[(1.4)] (No image; 8 minutes for questions 4--6) If the age of the Moon
  is 27.5 days, in which direction can we observe it and at what time?
  Select only one choice.
  \begin{description}
  \item[A.] Before sunrise to the east.
  \item[B.] After sunset to the west.
  \item[C.] Close to zenith at sunrise.
  \item[D.] At 9 am rising in the East.
  \end{description}
  \hfill[2.0 pt]
\item[(1.5)] Which of the Right Ascensions ($RA$) listed below will
  approximately coincide with the local meridian at Pune, India
  ($74^\circ$E) on April 10 at 19:15 local time?
  \begin{description}
  \item[A.] 6 h
  \item[B.] 7 h
  \item[C.] 8 h
  \item[D.] 9 h
  \item[E.] 10 h
  \end{description}
  \hfill[5.0 pt]
\item[(1.6)] At what time during this night (April 10) would we find Regulus
  [$\alpha$ Leo, $RA = 10^{\mathrm{h}}08^{\mathrm{m}}22.3^{\mathrm{s}}$,
  $Dec = +11^\circ 58' 02''$] at its maximum altitude as seen from Pune,
  India? Select only one choice.
  \begin{description}
  \item[A.] 18:00h
  \item[B.] 19:00h
  \item[C.] 20:00h
  \item[D.] 21:00h
  \item[E.] 22:00h
  \end{description}
  \hfill[4.0 pt]
\end{description}

For questions 7--10 the projected sky (Image 4) corresponds to latitude
$-30^\circ$ (South) and longitude $-74^\circ$ (West).

\begin{description}
\item[(1.7)] (Image 4, 2 minutes) Select the constellation pack that crosses
  the celestial equator on the image.
  \begin{description}
  \item[A.] Ori, Tau, Cet, Eri, Psc, Aqr, Aql, Ser
  \item[B.] Mic, Tau, Per, Cae, Mon, Aql, And, Ser
  \item[C.] Gru, Tau, Cha, Ara, Aqr, Aql, And, Aps
  \item[D.] For, Tri, Cet, Psc, Sct, Aql, And, Ser
  \end{description}
  \hfill[3.0 pt]
\item[(1.8)] (Image 4, 2 minutes) Choose the name of the stars indicated by
  the arrows (tick the correct cell).
  [Answer table: for star numbers 1, 2, 3, tick one of --- Achernar,
  Alnair, Altair, Alnath, Ankaa, Canopus.] \hfill[6.0 pt]
\item[(1.9)] (Image 4, 3 minutes) Are the following Messier objects present
  in the projected sky? Select YES or NO for each:
  M15 (Peg), M16 (Aql), M27 (Vul), M6 (Sco), M15 (Cas), M27 (CMi).
  \hfill[6.0 pt]
\item[(1.10)] (Image 4, 3 minutes) Choose the name of the selected stars
  (tick the correct cell). If the same object is listed more than once,
  choose only the option with the correct constellation.
  [Answer table: for star numbers 1, 2, 3, tick one of --- Miaplacidus
  (Car), Miaplacidus (Lac), Sham (Sge), Sham (For), The Persian (Gru), The
  Persian (Ind).] \hfill[3.0 pt]
\item[(1.11)] (Image 5, for questions 11 and 12, 4 minutes) From the list
  below select the option that contains the constellations that are along
  the line of the ecliptic:
  \begin{description}
  \item[A.] Psc, Cap, Sgr, Oph
  \item[B.] Lib, Sco, Oph, Sgr
  \item[C.] Cap, Sge, Lib, Vir
  \item[D.] Sgr, Oph, Sco, Vir
  \item[E.] CrA, Sco, Oph, Leo
  \end{description}
  \hfill[4.0 pt]
\item[(1.12)] Which constellation or part of constellation is encircled in
  the image?
  \begin{description}
  \item[A.] Sagittarius (Sgr)
  \item[B.] Corvus (Crv)
  \item[C.] Telescopium (Tel)
  \item[D.] Corona Australis (CrA)
  \end{description}
  \hfill[3.0 pt]
\item[(1.13)] (Image 6, 2 minutes) You are making an observation through a
  telescope pointed towards the centre of the galaxy. You find an
  interesting object as shown in image 6. What is this object?
  \begin{description}
  \item[A.] M42 -- Orion Nebula
  \item[B.] M31 -- Andromeda Galaxy
  \item[C.] NGC 2024 -- Flame Nebula
  \item[D.] M8 -- Lagoon Nebula
  \item[E.] M20 -- Trifid Nebula
  \end{description}
  \hfill[3.0 pt]
\item[(1.14)] (Image 7, 2 minutes) In the Messier Astronomical Catalogue
  there is a mysterious object: a galaxy. This object in the image cannot be
  easily identified. Some astronomers had said that it was a duplicate of
  another Messier object, but here we have it as an original and awesome
  galaxy. Which object is this?
  \begin{description}
  \item[A.] M31 -- Andromeda Galaxy
  \item[B.] M33 -- Triangulum Galaxy
  \item[C.] M51 -- Whirlpool Galaxy
  \item[D.] M101 -- Pinwheel Galaxy
  \item[E.] M102 -- Spindle Galaxy
  \end{description}
  \hfill[5.0 pt]
\item[(1.15)] (Image 8, 2 minutes) Estimate the magnitude of Shaula in the
  constellation Scorpius.
  \begin{description}
  \item[A.] 1
  \item[B.] 1.5
  \item[C.] 2
  \item[D.] 2.5
  \item[E.] 3
  \end{description}
  % note: the last option is misprinted "C." in the original paper
  \hfill[2.0 pt]
\item[(1.16)] (Image 9, 4 minutes) Several stars have been indicated in
  image 9. Tick the stars indicated in the table below.
  [Answer table: for star numbers 1, 2, 3, 4, tick one of --- Acubens,
  Adhara, Aludra, Alzirr, Arneb, Gomeisa, $\alpha$ Mon, Mebsuta, Mirzam,
  Wasat.] \hfill[8.0 pt]
\item[(1.17)] (Image 10, 4 minutes) A Messier object has been indicated in
  image 10. Tick the correct cell below to indicate what it is.
  [Answer table: for object numbers 1, 2, 3, 4, tick one of --- M41, M42,
  M46, M45, M48, M50, M64, M93.] \hfill[8.0 pt]
\end{description}
